% !PRINTED 620
\weberpage{620}
%<<CONT
belongs, must also be contained in $C$. $C$ is therefore a normal
divisor of $P$, and consequently, by Theorem I., the group $P$
must be linear.

VII.~\gesperrt{An irreducible equation of prime degree,
in which all the roots are rationally expressible in terms of
any two of them, is metacyclic.}

From these theorems we further infer a remarkable property
of metacyclic equations, first noticed by \gesperrt{Kronecker},
for the case of a real domain of rationality\footnote{Ueber algebraisch auflösbare Gleichungen. Monatsbericht der Berliner
Akademie, 14. April 1856.}.
If in such an equation two roots are real, then it follows
from VI that all its roots are real. But such an equation,
since it is of odd degree, must always have one real
root. Hence follows

VIII.~\gesperrt{An irreducible metacyclic equation of
odd prime degree with real coefficients has either
all its roots real or only one.}

If all the roots are real, the discriminant, as a product
of nothing but squares of real quantities $(x_h - x_k)^2$, is positive. If only
one root is real, then to every complex factor of the
discriminant $(x_h - x_k)^2$ there corresponds a conjugate one, and their product is
positive. Only when $x_h$ and $x_k$ are conjugate imaginary, so that $x_h - x_k$
is purely imaginary, is $(x_h - x_k)^2$ negative, and the discriminant
thus receives the factor $-1$ for every pair of conjugate imaginary
roots. The sign of the discriminant is therefore $(-1)^{\frac{n-1}{2}}$.
From this follows:

IX.~\gesperrt{When $n \equiv 1\ (\mathrm{mod}\ 4)$, the
discriminant is always positive; but if $n \equiv 3\ (\mathrm{mod}\ 4)$,
the sign of the discriminant decides which of the two
cases of Theorem VIII. occurs.}
